\chapter{Выводы}
\label{ch:conclusions}

Скорости откачки через капилляр после введения в систему течи $W = (27 \pm 10) \ \text{см}^3 / \text{с}$ оказалась меньше значения полученного по формуле (1) из зависимости P(t) в процессе откачки высоковакуумного баллона диффузионным насосом $W = (3.17 \pm 0.13) \cdot 10^2~\text{см}^3/\text{с}$. Это говорит о том, что после создания течи эффективная пропускная способность системы вместе с капилляром снизилась. Это не позволило достичь максимальной скорости откачки. Также это показывает, что скорость откачки зависит от пропускной способности трубок, потому что скорость откачки через узкий каппиляр оказалась меньше, чем через более широкую трубку.